\documentclass[11pt]{article}

\usepackage[margin=1in]{geometry}
\usepackage{amsmath, amssymb, amsthm}
\usepackage{booktabs}
\usepackage{graphicx}
\usepackage{hyperref}
\usepackage{mathtools}

\newtheorem{theorem}{Theorem}
\newtheorem{lemma}{Lemma}
\newtheorem{remark}{Remark}

\title{When and Why Does Value-Aware Model Learning Fail?\\
  \large A Diagnostic Suite for Objective Mismatch in Model-Based Reinforcement Learning}
\author{Firmansyah Sundana, Novanto Yudistira}
\date{Draft v1 (Phase 1d)}

\begin{document}
\maketitle

\begin{center}
\textbf{Status:} Draft v1 (Phase 1d). Tabular (Exp~1/2/1b) and continuous neural
(Exp~4) results included. Deep-baseline benchmark (Phase~2) pending provisioning.
\end{center}

\begin{abstract}
Model-based reinforcement learning fits a dynamics model to predict the next state,
typically by maximum likelihood, and then uses that model to improve a policy. The
objective mismatch problem~\cite{lambert2020objective} is that the model's likelihood
need not correlate with its ability to guide the policy. Value-aware model learning
(VaGraM~\cite{voelcker2022vagram}) reweights the model loss by the value function, but
is empirically known to sometimes underperform MLE. We isolate \emph{when and why} this
happens. We contribute (i) a theory of the value-aware weight estimator $w$ --- a
prediction-risk dominance result showing weighting can never improve the Bellman-target
prediction (it only adds a variance penalty $\sigma_w^2\,\mathbb{E}[V^2]/n$ and a bias
term), and a decision-risk crossover criterion showing weighting helps exactly when the
value-relevance signal exceeds the weight-noise penalty; (ii) a synthetic benchmark,
Distractor-Gym, that provably produces the failure regimes (spurious state dimensions
and sparse rewards); and (iii) two quantitative diagnostics --- policy-gradient
alignment (cosine between true and model-induced gradients) and the
$\lvert\delta_{\mathrm{TD}}\rvert \sim \lVert \nabla V\rVert \cdot \varepsilon_{\mathrm{model}}$
decomposition --- validated in a ground-truth tabular suite and transferred to a
continuous neural setting. We find that MLE misguides the policy exactly when data
coverage is concentrated where the value changes (alignment near zero or negative),
that value-aware weighting restores alignment there, and that value-aware weighting
collapses when its weights are degenerate: sparse rewards (noisy knife-edge value
gradient) or spread coverage. The decomposition holds with $R^2 = 0.99$ (dense) and
$0.83$ (sparse reward).
\end{abstract}

\section{Introduction}

Model-based RL (MBRL) is a sample-efficient framework for continuous
control~\cite{deisenroth2011pilco,janner2019mbpo,hafner2023dreamerv3,hansen2024tdmpc2}.
A dynamics model is fit to predict one-step transitions and then used by a policy or
planner. The standard objective for the model is maximum likelihood (MLE), which
minimizes prediction error uniformly over the state space. Lambert
et~al.~\cite{lambert2020objective} showed this objective can be \emph{mismatched} with
the downstream goal: the one-step likelihood is not always correlated with control
performance.

A principled family of remedies --- value-aware model learning
(VAML~\cite{farahmand2017vaml,farahmand2018itervaml}; VaGraM~\cite{voelcker2022vagram})
--- reweights the model loss by the value function, so that the model is accurate
\emph{where value changes}. VaGraM showed these losses help under small model capacity
and distracting state dimensions, yet the literature also documents regimes in which
value-aware losses \emph{underperform} MLE~\cite{voelcker2022vagram,cvaml2025,romi2026}.
There is currently no isolated, controlled characterization of when and why value-aware
weighting fails, and no modern diagnostic suite ported to current baselines.

This paper provides both. We make the following contributions:

\begin{enumerate}
\item \textbf{Theory of the weight estimator.} The value-aware weight
  $w = w(V, \nabla V)$ is itself estimated from finite data and is a random variable
  with bias and variance. We prove (i) a \emph{prediction-risk dominance} result:
  weighting can never improve the conditional-mean Bellman-target prediction, adding a
  variance penalty $\sigma_w^2\,\mathbb{E}[V^2]/n$ and a bias term (Theorem~\ref{thm:pred});
  (ii) an \emph{effective sample size} reduction $\mathrm{ESS} \sim n/(1+\sigma_w^2)$
  (Theorem~\ref{thm:ess}); and (iii) a \emph{decision-risk crossover} criterion:
  weighting improves the policy-relevant error iff the value-relevance signal
  $S = \left(\nabla V \cdot \mathrm{Cov}(w, Y)\right)^2$ exceeds the weight-noise penalty
  (Theorem~\ref{thm:crossover}).
\item \textbf{Distractor-Gym.} A synthetic environment with spurious state dimensions
  (complex, reward-irrelevant dynamics) and a sparse-reward variant, whose regime knobs
  provably produce the failure conditions of the theory.
\item \textbf{Diagnostics.} Policy-gradient alignment and the
  $\lvert\delta_{\mathrm{TD}}\rvert \sim \lVert \nabla V\rVert \cdot \varepsilon_{\mathrm{model}}$
  decomposition, measured in a ground-truth tabular suite and a continuous neural setting.
\end{enumerate}

\textbf{Findings.}
(a)~MLE \emph{fails to guide the policy} precisely when data coverage is concentrated
where the value changes: the cosine between the true and the model-induced policy
gradient falls to $\sim 0$ or negative (e.g., $-0.08$ under sparse reward + goal-directed
coverage), while value-aware weighting restores it to $\sim 0.9$.
(b)~Value-aware weighting \emph{fails} exactly when its weights are degenerate --- sparse
reward (weight sits on a noisy knife-edge value gradient) or spread coverage --- and in
these regimes MLE is strictly better.
(c)~On the Bellman-target prediction risk, MLE dominates value-aware weighting in every
regime we tested, and the gap is dominated by weight-estimator noise (oracle weights are
much better than estimated weights) --- reproducing and explaining VaGraM's observation
that value-aware losses ``tend to be inferior in practice to MLE.''

\section{Background and Related Work}

\textbf{Objective mismatch.} Lambert et~al.~\cite{lambert2020objective} formalized that a
model trained by one-step MLE likelihood and used for control may not correlate with
performance; they proposed an exploratory sample reweighting $w(y) = c \exp(-d(y))$.
Wei et~al.~\cite{wei2024unified} survey the decision-aware remedies that followed.

\textbf{Value-aware model learning.} Farahmand et~al.~\cite{farahmand2017vaml,farahmand2018itervaml}
introduced VAML and IterVAML, replacing the model loss with the induced value error.
VaGraM~\cite{voelcker2022vagram} uses the value gradient as the weight, and explicitly
documents two failure modes of naive value-aware learning: (a)~the loss does not account
for exploration/coverage of the state space, and (b)~the value function is often
non-smooth, so unbounded value-aware losses diverge. CVAML~\cite{cvaml2025} shows
sample-based value-aware losses are uncalibrated for stochastic models; ROMI~\cite{romi2026}
adds adaptive MLE+value weighting for offline RL. MOBILE~\cite{sun2023mobile} uses
model-Bellman inconsistency for offline pessimism.

\textbf{Diagnostics.} Lambert et~al.~\cite{lambert2020objective} measured the correlation
between one-step likelihood and return. We modernize this: our diagnostics are per-run,
quantitative, and cover gradient alignment and the value-error decomposition, ported
(Phase~2) to MBPO, DreamerV3, TD-MPC2, and VaGraM.

\section{Theory: The Weight Estimator \texorpdfstring{$w$}{w}}
\label{sec:theory}

\textbf{Setup.} Model family $P_\theta(s' \mid s, a)$; weighted loss
\begin{equation}
  L_w(\theta) \;=\; \mathbb{E}_{\mathcal{D}}\!\left[
    w(s, a, s')\, \ell\bigl(P_\theta(s' \mid s, a)\bigr)
  \right],
  \label{eq:weighted-loss}
\end{equation}
with MLE the special case $w = 1$. Value-aware weights: $w = \lvert V(s') - V(s)\rvert$
(VAML-1), $w = \lVert \nabla_s V(s')\rVert$ (VaGraM),
$w = \exp\bigl((V(s') - V_{\max})/\tau\bigr)$ (Lambert). The weight is never known; it
is estimated from a data-dependent value estimate, so
$\hat{w} = w(\hat{V}, \nabla \hat{V})$ has bias and variance. Full statements and proofs
in \texttt{paper/notes/theory.md}.

\begin{lemma}[decomposition]
\label{lem:decomp}
Let $V$ be twice differentiable with Hessian $H_V$, and let $e = \hat{s}' - s'$ be the
model prediction error. Then, with $\phi = \angle(\nabla_s V(s'), e)$,
\begin{align}
  \delta_{\mathrm{TD}}(s,a) &:= V(s') - V(\hat{s}') 
    \;=\; - \nabla_s V(s')^{\top} e - \tfrac{1}{2} e^{\top} H_V(\xi)\, e, \\
  \lvert \delta_{\mathrm{TD}} \rvert
    &\le \lVert \nabla_s V(s')\rVert\, \lVert e\rVert\, \lvert\cos\phi\rvert
       + \tfrac{1}{2} \lVert H_V \rVert_2\, \lVert e\rVert^2.
\end{align}
\end{lemma}

The second term is the \textbf{curvature residual} that the first-order
$\lVert \nabla V\rVert \cdot \varepsilon_{\mathrm{model}}$ factorization ignores. This
decomposition is exactly what value-aware weighting targets: VaGraM weights by
$\lVert \nabla V\rVert$ alone, so it is mis-specified wherever the curvature term or the
$\cos\phi$ distribution dominates. \emph{Empirical check (Exp~2 Part A):}
$R^2 = 0.99$ (dense reward) and $0.83$ (sparse step-function value), where the curvature
residual grows.

\begin{theorem}[prediction-risk dominance]
\label{thm:pred}
Let $y_1, \dots, y_n$ be iid draws from the true transition at a fixed $(s,a)$, and
$w_i \ge 0$ non-negative weights with $\mathbb{E}[w] = 1$, $\mathrm{Var}(w) = \sigma_w^2$,
independent of $y$. Then, to first order in $1/n$,
\begin{align}
  \mathrm{Var}(\hat{\mu}_w) &= \mathrm{Var}(\hat{\mu}_{\mathrm{MLE}})
    + \frac{\sigma_w^2\, \mathbb{E}[V^2]}{n}, \\
  \mathbb{E}[\hat{\mu}_w] &= \mathbb{E}[V(y)] + \mathrm{Cov}(w, V(y)) + O(1/n),
\end{align}
so the expected squared value-prediction error satisfies
\begin{equation}
  \mathbb{E}\!\left[\bigl(\hat{\mu}_w - \mathbb{E}[V]\bigr)^2\right]
  = \mathbb{E}\!\left[\bigl(\hat{\mu}_{\mathrm{MLE}} - \mathbb{E}[V]\bigr)^2\right]
    + \frac{\sigma_w^2\, \mathbb{E}[V^2]}{n}
    + \mathrm{Cov}(w, V(y))^2 + O(n^{-3/2}).
\end{equation}
\end{theorem}

Value-aware weighting \textbf{cannot improve the conditional-mean (Bellman-target)
prediction}: it adds variance and bias. Both grow when the weight estimator is noisy
(sparse reward $\Rightarrow$ noisy $\hat{V}$ $\Rightarrow$ large $\sigma_w^2$) and when
the value is flat (distractor dims $\Rightarrow$ $\mathrm{Cov}(w, V) \approx 0$).
\emph{Empirical check (Exp~1b):} the Bellman risk of MLE is at or below that of every
weighted model across all regimes; the estimated-weight model is strictly worse than the
oracle-weight model.

\begin{theorem}[effective sample size]
\label{thm:ess}
For mean-normalized weights, the Kish effective sample size satisfies
\begin{equation}
  \mathrm{ESS}(w) \;=\; \frac{\left(\sum_i w_i\right)^2}{\sum_i w_i^2}
  \;\sim\; \frac{n}{1 + \sigma_w^2} \;\le\; n.
\end{equation}
\end{theorem}

High-variance weights shrink the effective sample size of the weighted loss; this is the
same variance inflation as Theorem~\ref{thm:pred} expressed as sample loss.

\begin{theorem}[decision-risk crossover criterion]
\label{thm:crossover}
The \emph{decision} risk is the value-relevant prediction error linearized along the
value gradient,
\begin{equation}
  R_{\mathrm{dec}}(\hat{P})
    = \mathbb{E}_{(s,a)}\!\left[
      \Bigl( \nabla_s V(s') \cdot 
        \bigl( \mathbb{E}_{\hat{P}}[Y \mid s,a] - \mathbb{E}_{P^*}[Y \mid s,a] \bigr)
      \Bigr)^2
    \right].
\end{equation}
Applying Theorem~\ref{thm:pred} to the vector estimator gives
\begin{equation}
  R_{\mathrm{dec}}(\hat{P}_w) - R_{\mathrm{dec}}(\hat{P}_{\mathrm{MLE}})
  = \underbrace{\left(\nabla V \cdot \mathrm{Cov}(w, Y)\right)^2}_{\text{value-relevance signal}}
    - \underbrace{\lVert \nabla V\rVert^2\, \sigma_w^2\,
       \frac{\mathbb{E}[\lVert Y\rVert^2]}{n}}_{\text{weight-noise penalty}}
    + \underbrace{\text{(weight-estimator bias term)}}_{\text{bias}}.
\end{equation}
\end{theorem}

The weighted model improves the decision risk iff the signal exceeds the penalty. This
yields the crossover conditions:
\begin{itemize}
\item \textbf{Distractor dims:} $\nabla_{s_d} V = 0$ $\Rightarrow$ signal $0$
  $\Rightarrow$ MLE optimal.
\item \textbf{Sparse reward / noisy $\hat{V}$:} $\sigma_w^2$ large (weights on a noisy
  value cliff) $\Rightarrow$ penalty dominates $\Rightarrow$ MLE wins.
\item \textbf{Informative coverage with a strong gradient:} signal dominates
  $\Rightarrow$ value-aware wins.
\end{itemize}

The computable statistic $\mathrm{SNR}_w = \mathbb{E}[w\ell]^2 / \mathrm{Var}(w\ell)$
proxies signal over penalty. \emph{Empirical check (Exp~1b):} the alignment delta
$\cos(g_{\text{true}}, g_w) - \cos(g_{\text{true}}, g_{\mathrm{MLE}})$ is positive in
goal-coverage sparse regimes and negative precisely where $\mathrm{SNR}_w$ collapses; it
is a positive-but-noisy function of $\mathrm{SNR}_w$ (Pearson $\approx 0.25$--$0.34$),
motivating a sharper diagnostic (open item).

\section{Distractor-Gym}

A synthetic environment isolating the failure conditions. State $s = (s_c, s_d)$ with
$d_c$ control-relevant and $d_d$ distractor dimensions whose dynamics are complex
(linear, nonlinear chaotic, or stochastic white noise) and whose reward relevance is
zero, so $\nabla_{s_d} V = 0$ by construction.

\begin{itemize}
\item \textbf{Tabular suite} (ground-truth lab): 1D control axis $s_c \in [-3, 3]$ on a
  grid with actions $\{-1, 0, +1\}$ and quantized transition noise; distractor block
  evolves by a deterministic chaotic logistic map per dimension. Reward dense
  $r = \exp\!\left(-(x - \mathrm{goal})^2/2\right)$ or sparse
  $r = \mathbf{1}[\lvert x - \mathrm{goal}\rvert \le \mathrm{radius}]$. Exact value
  functions, value gradients, policy gradients, and Bellman risks are computable.
\item \textbf{Continuous suite:} Pendulum-v1 with an appended chaotic or white-noise
  distractor block (Gymnasium wrapper). Reward depends only on the control block.
\end{itemize}

Regime knobs swept: $d_d$ (distractor count), distractor dynamics class, reward sparsity
(via goal radius and discount $\gamma$), transition noise, coverage (uniform vs
goal-directed data collection), and model capacity. These produce the mismatch phase
diagram (MLE-win / value-aware-win / both-fail) that the diagnostics quantify.

\section{Experiment 1: Policy-Gradient Alignment}

\textbf{Setup.} For a fixed (deliberately off-center) softmax policy $\pi$, compute the
exact policy gradient under the true dynamics $g_{\text{true}}$ and under a model fitted
by each loss family $g_{\mathrm{model}}$; report
$\cos(g_{\text{true}}, g_{\mathrm{model}})$ (tabulated in Exp~1 over $d_d$, reward
sparsity, coverage; averaged over 5 seeds).

\textbf{Result.} Table~\ref{tab:exp1}. MLE \emph{misguides the policy} under sparse
reward + goal-directed coverage: alignment falls to $-0.08$ (its gradient points against
the true improvement direction), while value-aware VAML-1/Lambert restore alignment to
$\sim 0.9$. Under dense reward + uniform coverage, MLE is near-perfectly aligned
($1.00$) and value-aware weighting is equal or worse. Value-aware \emph{fails} in two
regimes: sparse reward + uniform coverage (VaGraM alignment $0.03$ vs MLE $0.99$ --- the
knife-edge weights concentrate the model on a razor-thin band) and dense + uniform
(VAML-1 drops to $0.88$). These are exactly the Theorem~\ref{thm:crossover} predictions.

\begin{table}[htbp]
\centering
\caption{Policy-gradient alignment $\cos(g_{\text{true}}, g_{\mathrm{model}})$ by regime
(averaged over 5 seeds).}
\label{tab:exp1}
\begin{tabular}{lcccc}
\toprule
Regime & MLE & VAML-1 & VaGraM & Lambert \\
\midrule
dense + uniform  & 1.00 & 0.88   & 1.00   & 0.31 \\
sparse + uniform & 0.99 & 0.95   & 0.03   & 0.95 \\
sparse + goal    & $-0.08$ & 0.89 & 0.29   & 0.91 \\
\bottomrule
\end{tabular}
\end{table}

\section{Experiment 2: Decomposition and Weight-Estimator Ablation}

\textbf{Part A (Lemma~\ref{lem:decomp}).} Across regimes, the factorization
$\lvert\delta_{\mathrm{TD}}\rvert \sim \lVert \nabla V\rVert \cdot \varepsilon_{\mathrm{model}} \cdot \lvert\cos\phi\rvert$
holds with slope $\sim 1.06$--$1.18$ and $R^2 = 0.99$ (dense) / $0.83$ (sparse). The
curvature residual ($0.006$ vs $0.022$) is the gap and grows with reward sparsity --- the
regime where gradient-only weighting is mis-specified.

\textbf{Part B (weight estimator vs policy effects).} Oracle weights (true $V$, $\nabla V$)
vs estimated weights ($\hat{V}$ from data). The estimated-weight model has consistently
higher Bellman risk than the oracle-weight model (e.g., VAML-1 risk $0.239$ vs $0.184$;
VaGraM $0.466$ vs $0.168$ at $d_d = 2$), isolating the weight-estimator noise term
$\sigma_w^2\, \mathbb{E}[V^2]/n$. Yet value-aware alignment improvements are robust to
this noise (align$_{\text{est}}$ and align$_{\text{ora}}$ both well above MLE in sparse
regimes). The value-aware benefit is decision-level (alignment) and does not come from
better Bellman prediction --- consistent with Theorems~\ref{thm:pred} and~\ref{thm:crossover}.

\section{Experiment 4: Continuous Neural Transfer}

Small MLP dynamics model on Pendulum-plus-distractors, MLE-MSE vs the VaGraM projection
loss $\bigl(\nabla r(s') \cdot (\hat{s}' - s')\bigr)^2$.

\begin{itemize}
\item \textbf{Pure projection ($\lambda = 1$) diverges} and fails to train a usable
  model --- reproducing VaGraM's documented divergence of naive value-aware losses.
\item \textbf{Bounded combination ($\lambda = 0.5$)} is stable and matches MLE on global
  MSE, while improving the decision-relevant metrics when unpredictable stochastic
  distractors consume capacity: value error $0.214$ vs $0.306$ ($d_d = 8$) and $0.524$
  vs $0.674$ ($d_d = 16$); directional alignment $0.971/0.942$ vs $0.952/0.905$. The
  benefit grows with distractor count and requires sufficient model capacity.
\end{itemize}

\section{Discussion}

\textbf{When MLE fails to guide the policy.} Under goal-concentrated coverage with sparse
reward, the MLE model's gradient misaligns ($\cos \sim 0$ or negative) because the value
cliff is thin and the model spreads its accuracy where the policy does not visit.
Value-aware weighting reallocates the model's effective mass to the cliff and restores
alignment.

\textbf{Why value-aware fails.} Two mechanisms, both predicted by the theory:
(1)~on the Bellman/prediction risk it \emph{cannot} help --- it only adds weight-noise
variance --- which reproduces VaGraM's ``inferior in practice'' observation;
(2)~on the decision risk it helps only when the value-relevance signal beats the
weight-noise penalty, which fails when weights are degenerate (sparse reward, spread
coverage).

\textbf{Practical rule of thumb.} Measure $\mathrm{SNR}_w$ and the gradient alignment
during training. If $\mathrm{SNR}_w$ is low or the alignment delta is non-positive,
value-aware weighting is not helping; the model is in the MLE-optimal regime.

\section{Limitations and Future Work}

\begin{itemize}
\item The $\mathrm{SNR}_w$ crossover statistic is a positive-but-noisy predictor; a
  sharper decision-aware statistic is open work.
\item Tabular results are exact but stylized; the continuous transfer used reward
  gradients as the value proxy, not a learned value function.
\item Phase~2: the diagnostics must be ported to MBPO, DreamerV3, TD-MPC2, and VaGraM
  (benchmark provisioning is an open item; \texttt{mbrl-lib} pins old \texttt{gym}).
\end{itemize}

\begin{thebibliography}{99}

\bibitem{deisenroth2011pilco}
M.~Deisenroth and C.~Rasmussen.
PILCO: A model-based and data-efficient approach to policy search.
In \emph{Proceedings of ICML}, 2011.

\bibitem{farahmand2017vaml}
A.~Farahmand, A.~Barron, and C.~Szepesv\'ari.
Value-aware model learning.
2017.

\bibitem{farahmand2018itervaml}
A.~Farahmand.
Iterative value-aware model learning.
In \emph{Advances in Neural Information Processing Systems}, 2018.

\bibitem{lambert2020objective}
N.~Lambert, B.~Amos, O.~Yadan, and R.~Calandra.
Objective mismatch in model-based reinforcement learning.
In \emph{Proceedings of L4DC}, 2020.

\bibitem{janner2019mbpo}
M.~Janner, J.~Fu, M.~Zhang, and S.~Levine.
When to trust your model: Model-based policy optimization.
In \emph{Advances in Neural Information Processing Systems}, 2019.

\bibitem{voelcker2022vagram}
C.~Voelcker, V.~Liao, A.~Garg, and A.~Farahmand.
Value gradient weighted model-based reinforcement learning (VaGraM).
In \emph{Proceedings of ICLR}, 2022.

\bibitem{sun2023mobile}
Y.~Sun, J.~Zhang, C.~Jia, H.~Lin, J.~Ye, and Y.~Yu.
Model-bellman inconsistency for model-based offline reinforcement learning (MOBILE).
In \emph{Proceedings of ICML}, 2023.

\bibitem{hafner2023dreamerv3}
D.~Hafner, J.~Pasukonis, J.~Ba, and T.~Lillicrap.
Mastering diverse domains through world models.
2023.

\bibitem{hansen2024tdmpc2}
N.~Hansen, H.~Su, and X.~Wang.
TD-MPC2: Scalable, robust world models for continuous control.
In \emph{Proceedings of ICML}, 2024.

\bibitem{wei2024unified}
R.~Wei, N.~Lambert, and A.~McDonald.
A unified view on solving objective mismatch in model-based reinforcement learning.
2024.

\bibitem{cvaml2025}
C.~Voelcker et~al.
Calibrated value-aware model learning (CVAML).
2025.

\bibitem{romi2026}
ROMI: Model-based offline reinforcement learning via robust value-aware model learning.
In \emph{Proceedings of ICLR}, 2026.

\end{thebibliography}

\end{document}